% !TeX root = ../../Book2.tex
% 纲要标题：### 19.1 Ore 局部化
% 纲要位置：计划纲要.md，第 1670 行。

\section{Ore 局部化}
\label{b2:sec:1901}

交换环中把分母移到乘积的一侧，依靠的是元素可以交换。非交换环里，写下 \(as^{-1}\) 之后，还须解释 \(s^{-1}b\) 如何重新写成分母在右侧的形式。Ore 条件正是为这一步提供可解的方程。本节只讨论由双侧非零因子组成的分母集，从而原环能够单射到分式环。\cref{b2:prop:zero-divisor-inversion}将说明，若把原环中的零因子变成可逆元，单射性就必然丢失。

% 纲要要点（供目录核对）：
%
% * 分母在非交换情形的障碍；由实际右分式乘积产生 Ore 方程，正文证明只逆一个正则元素也会失败、逆零因子必破坏单射性
% * 左、右 Ore 条件；按最终分母位置区分，共同分母属于分母集而正则乘子不必属于其中
% * 典型局部化及其泛性质；先构造加法分式类，再由左乘算子复合给出结合乘法，解释被迫可逆的乘子及非中心矩阵计算
%

\subsection{分母在非交换情形的障碍}
\label{b2:subsec:1901:01}

本节设 \(R\ne0\) 为含单位元的环。若 \(s\in R\) 满足
\[
sx=0\Longrightarrow x=0,\qquad xs=0\Longrightarrow x=0,
\]
称 \(s\) 为\term[zhengze-yuansu]{正则元素}[regular element]，即它既不是左零因子，也不是右零因子。正则元素的乘积仍正则，一也是正则元素。

\ProseParagraphBreak
假设 \(R\) 已嵌入某个环 \(Q\)，分母集 \(S\) 中的元素在 \(Q\) 中可逆，并且每个 \(Q\) 元素都能写为右分式。对 \(p,a\in R\) 及 \(s,r\in S\)，乘积为
\[
(ps^{-1})(ar^{-1})=p(s^{-1}a)r^{-1}.
\]
要使中间的 \(s^{-1}a\) 也成为右分式，必须将它写成 \(bt^{-1}\)，其中 \(b\in R,t\in S\)。这个等式两侧分别左乘 \(s\)、右乘 \(t\)，就得到
\[
at=sb.
\]
因此共同倍数方程正是乘法中移换分母的要求。取得这样的 \(b,t\) 后，乘积便能整理为 \((pb)(rt)^{-1}\)。

\ProseParagraphBreak
以自由代数 \(R=k\langle x,y\rangle\) 为例。由\cref{b2:prop:free-associative-algebra}，有限字构成基。对非零多项式按字长、同长度内再按字典序选最大字，乘积中字长达到最大时，两因子都须各自达到最大字长；在这些固定长度的字中，任一因子按字典序变小，连接后的字也变小。因此乘积的最大字只能由两个最大字连接而成，系数为两个非零首系数的乘积，故乘积非零。因此所有非零元素正则。

\ProseParagraphBreak
但 \(xR\cap yR=0\)，因为第一项中的非零字都以 \(x\) 开头，第二项中的非零字都以 \(y\) 开头。若 \(a=x,s=y\) 能满足上述关系且 \(t\ne0\)，就有 \(0\ne xt=yb\in xR\cap yR\)，矛盾。所以这个非交换整环不能通过“所有非零元素作右分母，且每个元素是一个右分式”的方式得到分式除环。这不排除其他类型的扩张环，只说明通常的单分式表达需要额外条件。即使把分母限制为一个正则元素 \(x\) 的各次幂，所需方程也未必可解，见\cref{b2:prop:single-element-ore-obstruction}。

\subsection{左、右 Ore 条件}
\label{b2:subsec:1901:02}

\begin{definition}\label{b2:def:ore-set}
设 \(R\ne0\) 为含幺环，\(S\subseteq R\) 含一、乘法封闭，且全部元素均双侧正则。对任意 \(a\in R,s\in S\)，若能取 \(t\in S,b\in R\) 使 \(at=sb\)，则称 \(S\) 满足\term[you-Ore-tiaojian]{右 Ore 条件}[right Ore condition]，并称这样的 \(S\) 为\term[zhengze-you-Ore-ji]{正则右 Ore 集}[regular right Ore set]；在本节的正则情形，也称右 Ore 分母集。若对任意 \(a\in R,s\in S\) 都能取 \(t\in S,b\in R\) 使 \(ta=bs\)，则称 \(S\) 为正则左 Ore 集。两个方程及其分式方向为
\[
\begin{aligned}
\text{右 Ore：}&\quad at=sb,\qquad as^{-1};\\
\text{左 Ore：}&\quad ta=bs,\qquad s^{-1}a.
\end{aligned}
\]
这里的左右按最终分式中分母的位置命名。左条件就是反环中的右条件；即使分母都正则，两条件也不自动相同。
\end{definition}

\begin{lemma}\label{b2:lem:ore-common-denominators}
设 \(R\ne0\) 为含幺环，\(S\subseteq R\) 为正则右 Ore 集。
\begin{enumerate}
\item 若 \(s,su\in S\)，则 \(u\) 正则，但不一定属于 \(S\)。
\item 任意 \(s,t\in S\) 有共同右倍数 \(su=tv\in S\)，其中乘子 \(u,v\) 均正则。任意有限多个分母也有这样的共同右倍数。
\end{enumerate}
\end{lemma}
\begin{proof}
若 \(ux=0\)，则 \((su)x=0\)，由 \(su\) 正则得 \(x=0\)。对分子 \(s\) 和分母 \(su\) 使用 Ore 条件，取 \(v\in S,b\in R\) 使 \(sv=(su)b\)。消去左侧正则因子 \(s\)，得 \(v=ub\)。若 \(xu=0\)，则 \(xv=xub=0\)，由 \(v\) 正则得 \(x=0\)。故 \(u\) 双侧正则。

对分子 \(t\)、分母 \(s\) 使用 Ore 条件，得 \(tv=su\)，其中 \(v\in S\)。共同值属于 \(S\)，第一部分保证 \(u\) 也正则。有限多个分母逐次与已经取得的共同分母再作此操作即可；每个旧分母都仍右整除新的共同值。
\end{proof}

共同分母必须属于 \(S\)，放大分母所用的乘子则只需正则，以便在比较分子时消去。例如在 \(\mathbb Z\) 中取 \(S=\{2^i6^j:i,j\ge0\}\)，有 \(2,6\in S\) 及 \(2\cdot3=6\)，但乘子三不属于 \(S\)。

\ProseParagraphBreak
在交换环中，或更一般地在分母都位于中心时，取 \(t=s\)、\(b=a\) 即满足右 Ore 条件，左条件也成立。一般情形的共同倍数则可能需要非平凡的乘子；不能直接把共同分母写成 \(st\)，因为 \(t\) 未必右整除 \(st\)。

\subsection{典型局部化及其泛性质}
\label{b2:subsec:1901:03}

共同右分母足以定义分式的加法，而非交换乘法还要处理分母与分子相遇时的次序。先构造加法群，再将分式实现成这个群上的左乘算子，就能用算子复合取得结合乘法。若 \(su=tv=w\in S\)，预期的换分母等式是 \(as^{-1}=(au)w^{-1}\)、\(bt^{-1}=(bv)w^{-1}\)；因此把两分式视为相同，正是要求它们通分后有相同的分子。

\begin{lemma}[右分式类的加法]\label{b2:lem:ore-additive-classes}
设 \(R\ne0\) 为含幺环，\(S\subseteq R\) 为正则右 Ore 集。在 \(R\times S\) 上规定
\begin{equation}\label{b2:eq:ore-equivalence}
(a,s)\sim(b,t)
\quad\Longleftrightarrow\quad
\exists u,v\in R:\ su=tv\in S,\quad au=bv.
\end{equation}
则 \(\sim\) 是等价关系。其商集 \(F=(R\times S)/{\sim}\) 上的公式
\[
[a,s]+[b,t]=[au+bv,w]\quad(su=tv=w\in S),
\qquad -[a,s]=[-a,s]
\]
给出交换群结构。其零元满足
\begin{equation}\label{b2:eq:ore-zero-numerator}
[a,s]=0\quad\Longleftrightarrow\quad a=0.
\end{equation}
\end{lemma}
\begin{proof}
自反性与对称性直接成立。证明传递性时，设 \(su=tv=w\)、\(au=bv\)，又有 \(tu'=rv'=w'\)、\(bu'=cv'\)。取共同右倍数 \(w\alpha=w'\beta\in S\)。从 \(tv\alpha=tu'\beta\) 消去左侧正则元 \(t\)，得到 \(v\alpha=u'\beta\)，于是
\[
s u\alpha=r v'\beta,\qquad
a u\alpha=bv\alpha=bu'\beta=cv'\beta.
\]
故 \((a,s)\sim(c,r)\)。

先核对一个换分母事实。若 \((a,s)\sim(b,t)\)，则对任意共同右分母 \(w=su=tv\in S\)，都有 \(au=bv\)。事实上，取等价关系中原有的共同分母 \(w_0=su_0=tv_0\)，再取 \(w_0\alpha=w\beta\in S\)。消去 \(s,t\) 得 \(u_0\alpha=u\beta\)、\(v_0\alpha=v\beta\)，故
\[
(au-bv)\beta=(au_0-bv_0)\alpha=0.
\]
由\cref{b2:lem:ore-common-denominators}，乘子 \(\beta\) 正则，因而 \(au=bv\)。

固定两对代表元，若分别选 \(w=su=tv\) 与 \(w'=su'=tv'\)，取 \(W=w\alpha=w'\beta\in S\)。消去 \(s,t\) 得 \(u\alpha=u'\beta\) 及 \(v\alpha=v'\beta\)，所以两个候选和扩张到 \(W\) 后分子相同。这证明共同分母的选择无关。

若把第一对代表元 \((a,s)\) 换为等价的 \((a',s')\)，而 \((b,t)\) 保持不变，分别取 \(w=su=tv\)、\(w'=s'u'=tv'\)，再取 \(W=w\alpha=w'\beta\in S\)。由刚证的换分母事实，\(a u\alpha=a'u'\beta\)；由 \(tv\alpha=tv'\beta\) 消去 \(t\)，得 \(v\alpha=v'\beta\)。因此
\[
(au+bv)\alpha=(a'u'+bv')\beta,
\]
两个候选和等价。若改换第二对代表元 \((b,t)\sim(b',t')\)，固定 \((a,s)\)，分别取 \(w=su=tv\)、\(w'=su'=t'v'\)，再取 \(W=w\alpha=w'\beta\in S\)。消去 \(s\) 得 \(u\alpha=u'\beta\)；换分母事实给出 \(bv\alpha=b'v'\beta\)。于是
\[
(au+bv)\alpha=(au'+b'v')\beta,
\]
两个候选和仍等价。依次换两对代表元即可处理任意代表元选择。

由\cref{b2:lem:ore-common-denominators}，三个类可取同一共同右分母，结合律与交换律于是归结为 \(R\) 的加法；\([0,1]\) 是零元，\([-a,s]\) 是 \([a,s]\) 的负元。最后，\([a,s]=[0,1]\) 时存在 \(u,v\in R\) 使 \(su=v\in S\)、\(au=0\)；由 \(u\) 正则得 \(a=0\)。反向由共同分母 \(s\) 直接得到。
\end{proof}

\begin{theorem}[正则分母的右 Ore 局部化]\label{b2:thm:right-ore-localization}
设 \(R\ne0\) 为含幺环，\(S\subseteq R\) 为正则右 Ore 集，则存在含幺环 \(Q=RS^{-1}\) 及保幺单射 \(R\rightarrow Q\)，使 \(S\) 中每个元素在 \(Q\) 中可逆，且每个 \(Q\) 元素都可写为 \(as^{-1}\)，其中 \(a\in R,s\in S\)。对任意保幺环同态 \(f:R\rightarrow B\)，若全部 \(f(s)\) 在 \(B\) 中可逆，则存在唯一保幺环同态
\[
\widetilde f:Q\longrightarrow B,\qquad
\widetilde f(as^{-1})=f(a)f(s)^{-1}.
\]
\end{theorem}
\begin{proof}
由\cref{b2:lem:ore-additive-classes}得到右分式类的加法群 \(F\)，在它的自同态环 \(\operatorname{End}_{\mathbb Z}(F)\) 中构造分式环。

对 \(a\in R\) 定义加法群自同态 \(L_a[b,t]=[ab,t]\)。等价关系和加法公式说明它良定义且可加，且 \(L_{a+b}=L_a+L_b\)、\(L_aL_b=L_{ab}\)、\(L_1=\operatorname{id}\)。对分母 \(s\in S\)，还须证明 \(L_s\) 是自同构。它单射，因为 \([sb,t]=0\) 迫使 \(sb=0\)，进而 \(b=0\)。它也满射：对 \([b,t]\)，由 Ore 条件取 \(v\in S,u\in R\) 使 \(bv=su\)。因 \(tv\in S\)，可以通分为
\[
[b,t]=[bv,tv]=L_s[u,tv].
\]
故 \(L_s\) 可逆。

在 \(\operatorname{End}_{\mathbb Z}(F)\) 中考虑所有算子 \(L_aL_s^{-1}\)。若 \(su=tv=w\in S\)，由 \(L_w=L_sL_u=L_tL_v\) 可知 \(L_u=L_s^{-1}L_w\)、\(L_v=L_t^{-1}L_w\) 也可逆。这一步用的是 \(s,t,w\) 都已被逆转，而不只是乘子 \(u,v\) 正则。于是
\[
L_aL_s^{-1}=L_{au}L_w^{-1},\qquad
L_bL_t^{-1}=L_{bv}L_w^{-1}.
\]
因此它们对加法与负元封闭，且等价分式给出相同算子。若 \(bv=su\)、\(v\in S\)，则
\[
L_s^{-1}L_b=L_uL_v^{-1},
\]
从而
\begin{equation}\label{b2:eq:right-ore-product}
\begin{aligned}
(L_aL_s^{-1})(L_bL_t^{-1})
&=L_{au}L_v^{-1}L_t^{-1}\\
&=L_{au}L_{tv}^{-1}.
\end{aligned}
\end{equation}
所以这些算子组成含单位元的子环。它们与 \(F\) 的元素一一对应：\(L_aL_s^{-1}\) 在 \([1,1]\) 处的值恰为 \([a,s]\)，因为 \(L_s[1,s]=[1,1]\)。故相同算子必来自相同分式类，反向已由共同分母证明。将这个子环的乘法转移到 \(F\)，得到结合环 \(Q\)。乘法结合律由算子复合继承，分配律也不需另作分母猜测。

映射 \(a\mapsto L_a\) 为保单位元环同态，且由在 \([1,1]\) 处取值及式\eqref{b2:eq:ore-zero-numerator}可知它单射。每个 \(L_s^{-1}=L_1L_s^{-1}\) 也在 \(Q\) 中，所以分母可逆。于是将 \([a,s]\) 记作 \(as^{-1}\) 后，式\eqref{b2:eq:right-ore-product}给出具体乘法
\[
(as^{-1})(bt^{-1})=(au)(tv)^{-1},
\qquad bv=su,\quad v\in S.
\]

最后核对泛性质。若 \(su=tv=w\in S\)，则 \(f(u)=f(s)^{-1}f(w)\) 可逆：即使 \(u\notin S\)，任何把 \(s,w\) 变成单位的同态也会迫使 \(u\) 变成单位。并有
\[
f(a)f(s)^{-1}=f(au)f(w)^{-1}.
\]
因此式\eqref{b2:eq:ore-equivalence}的相同分子保证 \(\widetilde f\) 良定义，共同分母公式保证可加。由 \(bv=su\)，得到 \(f(s)^{-1}f(b)=f(u)f(v)^{-1}\)，恰好验证乘法公式被保持。单位也被保持。每个元素都是 \(as^{-1}\)，其像已由 \(f\) 及分母求逆决定，故延拓唯一。
\end{proof}
\symidx{\(RS^{-1}\)：右 Ore 分式环}

这称为\term[Ore-jubuhua]{Ore 局部化}[Ore localization]。对反环应用同一定理，得到左 Ore 分式环 \(S^{-1}R\)。若同一个 \(S\) 同时满足左右条件，两种环都满足使 \(S\) 中元素可逆的同一个泛性质。分别延拓原环到另一分式环的映射，得到两个方向的同态；其复合都在 \(R\) 上为恒等，泛性质中的唯一性便使复合在整个分式环上也为恒等。因此两种构造具有保持 \(R\) 的互逆同构。这里所用的唯一性论证与\cref{b2:prop:universal-object-uniqueness}中积、余积的情形相同。

\begin{proposition}\label{b2:prop:ore-clearing-denominators}
设 \(R\ne0\) 为含幺环，\(S\subseteq R\) 为正则右 Ore 集，\(Q=RS^{-1}\) 为其右 Ore 局部化。对任意有限组 \(q_1,\ldots,q_m\in Q\)，存在共同右分母 \(s\in S\) 使全部 \(q_i s\in R\)。若 \(I\) 为 \(R\) 的右理想，则
\[
IQ=\{\,as^{-1}:a\in I,\ s\in S\,\}.
\]
每个非零右 \(R\) 子模 \(X\subseteq Q\) 都与 \(R\) 有非零交。
\end{proposition}
\begin{proof}
先把各分母取共同右倍数，分子随之扩张，就得到第一项。\(IQ\) 中元素为有限和 \(\sum_i a_iq_i\)、\(a_i\in I\)，通分后分子仍属于右理想 \(I\)，得到第二式的一向；另一向由 \(as^{-1}\in IQ\) 直接成立。最后，若 \(0\ne as^{-1}\in X\)，右乘 \(s\in R\) 得到 \(0\ne a\in X\cap R\)，其中非零性来自分母在 \(Q\) 中可逆。
\end{proof}

例如令 \(C\) 为交换整环，\(R=M_n(C)\)，分母取非零标量矩阵 \(sI_n\)。它们在中心且正则，局部化得到
\[
RS^{-1}\cong M_n(\operatorname{Frac}C).
\]
\symidx{\(\operatorname{Frac}C\)：交换整环的分式域}
每个右分式逐条目取标量分母，反向对一个分式域矩阵的有限多个条目取共同分母，便写成一个矩阵右分式；两种操作互逆且保持加乘。原环可能有许多矩阵零因子，但这些元素没有被放入分母集。

\ProseParagraphBreak
也可以让分母不在中心，此时 \(s^{-1}a\) 与 \(as^{-1}\) 可能不同，换分母还须改变分子。取 \(R=M_2(\mathbb Z)\)，令 \(S\) 为其中行列式非零的全部矩阵。它们在 \(M_2(\mathbb Q)\) 中可逆，因此在 \(R\) 中双侧正则，且 \(S\) 含一、乘法封闭。给定 \(a\in R,s\in S\)，选正整数 \(q\) 清除 \(s^{-1}a\) 的全部条目分母，令 \(t=qI_2\in S\)、\(b=s^{-1}at\in R\)，便有 \(at=sb\)。所以 \(S\) 满足右 Ore 条件。所得分式环仍为 \(M_2(\mathbb Q)\)：每个右分式可在该矩阵环中解释，而每个有理矩阵都能用非零整数标量矩阵作共同分母；若右分式在其中为零，其分子也必为零。

\ProseParagraphBreak
具体取
\[
s=\begin{pmatrix}1&1\\0&2\end{pmatrix},\qquad
a=E_{21},\qquad t=2I_2,\qquad
b=\begin{pmatrix}-1&0\\1&0\end{pmatrix}.
\]
直接相乘得到
\[
at=\begin{pmatrix}0&0\\2&0\end{pmatrix}=sb.
\]
于是非中心分母可以按 Ore 方程移到右侧，但分子也随之改变：
\[
s^{-1}a=bt^{-1}
=\begin{pmatrix}-\frac12&0\\\frac12&0\end{pmatrix},
\qquad
as^{-1}=\begin{pmatrix}0&0\\1&-\frac12\end{pmatrix}.
\]
例如两个右分式的乘积为
\[
\begin{aligned}
(E_{11}s^{-1})(as^{-1})
&=E_{11}b\,t^{-1}s^{-1}\\
&=E_{11}b(st)^{-1}\\
&=-E_{11}(2s)^{-1}
=\begin{pmatrix}-\frac12&\frac14\\0&0\end{pmatrix}.
\end{aligned}
\]
这里先用 \(s^{-1}a=bt^{-1}\) 处理两个分式之间相遇的分母与分子，再用 \(t^{-1}s^{-1}=(st)^{-1}\) 合并分母；次序来自实际乘法。

\begin{proposition}\label{b2:prop:single-element-ore-obstruction}
设 \(k\) 为域，\(R=k\langle x,y\rangle\) 为两个生成元上的自由结合含幺 \(k\) 代数，\(S=\{x^n:n\ge0\}\)。则 \(S\) 含一、乘法封闭，全部元素双侧正则，但不满足右 Ore 条件。
\end{proposition}
\begin{proof}
由\cref{b2:prop:free-associative-algebra}，\(x,y\) 上的有限字是 \(R\) 的一组 \(k\) 基，空字表示一。对每个 \(n\ge0\)，在字的左侧或右侧连接 \(x^n\)，都把不同的字送到不同的字。因此左乘与右乘 \(x^n\) 都是单射，\(x^n\) 双侧正则。\(x^0=1\) 及 \(x^i x^j=x^{i+j}\) 还说明 \(S\) 含一、乘法封闭。

若满足右 Ore 条件，对 \(a=y,s=x\) 就能取 \(t=x^n\in S\)、\(b\in R\)，使 \(yx^n=xb\)。左侧是以 \(y\) 开头的非零基字，右侧的每个非零基字都以 \(x\) 开头。字基的线性无关性排除该等式，所以右 Ore 条件失败。
\end{proof}

\begin{proposition}\label{b2:prop:zero-divisor-inversion}
设 \(R\ne0\)、\(B\) 为含幺环，\(f:R\to B\) 为保幺单射。若 \(s\in R\) 的像 \(f(s)\) 可逆，则 \(s\) 双侧正则。

具体地，设 \(k\) 为域，\(R=k\times k\)、\(e=(1,0)\)。第一坐标投影 \(\pi:R\to k\) 把 \(e\) 送到单位元，且对每个含幺环 \(B\) 及每个使 \(f(e)\) 可逆的保幺环同态 \(f:R\to B\)，都存在唯一保幺环同态 \(\bar f:k\to B\) 使 \(f=\bar f\pi\)。但是 \(\ker\pi=0\times k\ne0\)，任何这样的 \(f\) 都不是单射。
\end{proposition}
\begin{proof}
若 \(sx=0\)，则 \(f(s)f(x)=0\)，乘以 \(f(s)^{-1}\) 得 \(f(x)=0\)，再由单射性得 \(x=0\)。若 \(xs=0\)，在右侧乘以 \(f(s)^{-1}\) 同样得到 \(x=0\)。所以可嵌入后变成单位的元素必须双侧正则。

在第二个例子中，\(e^2=e\)。可逆的 \(f(e)\) 满足 \(f(e)^2=f(e)\)，消去一个因子得到 \(f(e)=1_B\)，于是 \(f(1-e)=0\)。每个第二坐标元素都为 \((0,b)=(1-e)(0,b)\)，故被 \(f\) 送到零。定义 \(\bar f(a)=f(a,0)\)，它保持加法与乘法，并把一送到 \(f(e)=1_B\)，因而为保幺环同态。又有 \(f(a,b)=f(a,0)=\bar f(\pi(a,b))\)，所以 \(f=\bar f\pi\)。投影 \(\pi\) 满射，保证这种分解唯一；它本身满足 \(\pi(e)=1\)，且核恰为非零的第二坐标因子。因此使 \(e\) 可逆必然零化原环中的非零元素。
\end{proof}
